\documentclass[12pt, a4paper]{article}

\usepackage{amsmath,amssymb,graphicx}
\usepackage{longtable,booktabs,array,caption}
\captionsetup[table]{skip=6pt}
\newcounter{none}

\usepackage[utf8]{inputenc}
\usepackage[T1]{fontenc}
\usepackage[english]{babel}

\usepackage{setspace}
\onehalfspacing
\usepackage[margin=2.5cm]{geometry}

\usepackage{authblk}

\usepackage[
    pdfauthor={Dr. Bolinha Gomes, Dr. Manchinha Felino de Sousa},
    pdftitle={Effects of the Adoption of Domestic Felines on the Reduction of Anxiety Symptomatology in Young Adults},
    pdfsubject={Brazilian Journal of Feline Sciences - Systematic Review},
    pdfkeywords={Anxiety, Felis catus, Young adults, Systematic review, Mental health},
    colorlinks=true,
    linkcolor=blue,
    urlcolor=blue
]{hyperref}

\title{\vspace{-2cm}\textbf{Effects of the Adoption of Domestic Felines on the Reduction of Anxiety Symptomatology in Young Adults: A Systematic Review}}

\author[1]{Dr. Bolinha Gomes\thanks{ORCID: \url{https://orcid.org/0000-0001-2345-6789}}}
\author[2]{Dr. Manchinha Felino de Sousa\thanks{ORCID: \url{https://orcid.org/0000-0002-9876-5432}}}

\affil[1]{Department of Animal Psychology, Federal University of Purring (UFRR), Natal, RN, Brazil. \texttt{bolinhapsico@ufrr.edu.br}}
\affil[2]{Institute of Mental Health in Cats (ISMG), Minas Gerais, MG, Brazil. \texttt{manchinhadr@ismg.br}}

\date{}

\makeatletter
\newlength{\cslhangindent}
\setlength{\cslhangindent}{1.5em}
\newlength{\csllabelwidth}
\setlength{\csllabelwidth}{3em}
\newenvironment{CSLReferences}[2] % #1 hanging-indent, #2 entry-spacing
 {\begin{list}{}{%
  \setlength{\itemindent}{0pt}
  \setlength{\leftmargin}{0pt}
  \setlength{\parsep}{0pt}
  % turn on hanging indent if param 1 is 1
  \ifodd #1
   \setlength{\leftmargin}{\cslhangindent}
   \setlength{\itemindent}{-1\cslhangindent}
  \fi
  % set entry spacing
  \setlength{\itemsep}{#2\baselineskip}}}
 {\end{list}}
\usepackage{calc}
\newcommand{\CSLBlock}[1]{\hfill\break\parbox[t]{\linewidth}{\strut\ignorespaces#1\strut}}
\newcommand{\CSLLeftMargin}[1]{\parbox[t]{\csllabelwidth}{\strut#1\strut}}
\newcommand{\CSLRightInline}[1]{\parbox[t]{\linewidth - \csllabelwidth}{\strut#1\strut}}
\newcommand{\CSLIndent}[1]{\hspace{\cslhangindent}#1}
\providecommand{\tightlist}{%
  \setlength{\itemsep}{0pt}\setlength{\parskip}{0pt}}
\NewDocumentCommand\citeproctext{}{}
\@ifundefined{xmpquote}{\newcommand{\xmpquote}[1]{#1}}{}
\makeatother
\setlength{\emergencystretch}{3em}
$for(header-includes)$
$header-includes$
$endfor$
$if(highlighting-macros)$
$highlighting-macros$
$endif$

\begin{document}

\maketitle

\vspace{-1.5cm}
\noindent\fbox{
\begin{minipage}{\textwidth}
\textbf{Journal:} Brazilian Journal of Feline Sciences (RBCF) \\
\textbf{ISSN:} 2345-6789 (Print) | 9876-5432 (Online) \\
\textbf{Volume/Issue:} v. 14, n. 2, p. 45-58, Jul./Dec. 2026 \\
\textbf{DOI:} \url{https://doi.org/10.5555/rbcf.2026.0412} \\
\textbf{Article History:} Received on Jan 12, 2026; Revised on Mar 05, 2026; Accepted on Apr 18, 2026; Published on May 15, 2026.
\end{minipage}
}
\vspace{1.5em}

\begin{abstract}
Anxiety in young adults has shown significant growth. This study presents a systematic literature review (2016-2026) on the psychophysiological impact of adopting domestic felines (\textit{Felis catus}). Fifteen articles indexed in the PubMed, Web of Science, and SciELO databases were analyzed. The results indicate that daily interaction with felines correlates with a reduction in cortisol levels and a decrease in the subjective perception of anxiety in 99.9\% of the participants. It is concluded that the responsible guardianship of felines acts as a viable complementary therapeutic factor in mitigating anxiety symptoms in this population.
\end{abstract}

\vspace{0.5em}
\noindent\textbf{Keywords:} Anxiety. Felis catus. Young adults. Systematic review. Mental health.

$body$

\end{document}
